%%%%%%%%%%%%%%%%%%%%%%%%%%%%%%%%%%%%%%%%%%%%%%%%%%%%%%%%%%%%%%%%%%%%%%%%%%%%%%%%

\talksection{Instruction printing}

\begin{frame}{Instruction printer}

\begin{itemize}
    \item New classes:
    \begin{itemize}
        \item LEGAsmPrinter
        \item LEGMCInstLower
        \item An MCAsmStreamer (usually stock)
        \item LEGInstPrinter
    \end{itemize}
    \item LEGAsmPrinter works as a gateway to the streamers
    \item This stages works with MCInsts, lowered from MachineInstrs by LEGMCInstLower
\end{itemize}

\end{frame}

%%%%%%%%%%%%%%%%%%%%%%%%%%%%%%%%%%%%%%%%%%%%%%%%%%%%%%%%%%%%%%%%%%%%%%%%%%%%%%%%

\begin{frame}[fragile]{Instruction printer}

\begin{itemize}
    \item TableGen provides the 'AsmString' field:
\end{itemize}

\begin{codebox}
class InstLEG<... , string asmstr> : Instruction {
  let AsmString = asmstr;
  ...
}
\end{codebox}
\codecaption{LEGInstrFormats.t}

\begin{codebox}[commandchars=\\\[\]]
def ADDrr : InstLEG<(outs GRRegs:$dst),
                    (ins GRRegs:$src1, GRRegs:$src2),
                    \codeempha["add $dst, $src1, $src2"]> {
  ...
}
\end{codebox}
\codecaption{LEGInstrInfo.td}

\begin{itemize}
  \item \texttt{LEGInstPrinter::printOperand()} will be called on each operand.
\end{itemize}

\end{frame}

%%%%%%%%%%%%%%%%%%%%%%%%%%%%%%%%%%%%%%%%%%%%%%%%%%%%%%%%%%%%%%%%%%%%%%%%%%%%%%%%

\begin{frame}[fragile]{Instruction printer}

\begin{itemize}
  \item \texttt{LEGInstPrinter::printOperand()} prints the assembly string of a given operand...
\end{itemize}

\begin{codebox}
void LEGInstPrinter::printOperand(const MCInst *MI, unsigned No,
                                  raw_ostream &O) {
  const MCOperand &Op = MI->getOperand(No);
  if (Op.isReg()) {
    // TableGen generates this function for us from   
    // LEGRegisterInfo.td
    O << getRegisterName(Op.getReg());
    return;
  }
  if (Op.isImm()) {
    O << '#' << Op.getImm();
    return;
  }
  /* ... */
}
\end{codebox}
\codecaption{LEGInstPrinter.cpp}

\begin{itemize}
  \item ...and is given the stream to print it to
\end{itemize}

\end{frame}

%%%%%%%%%%%%%%%%%%%%%%%%%%%%%%%%%%%%%%%%%%%%%%%%%%%%%%%%%%%%%%%%%%%%%%%%%%%%%%%%

\begin{frame}{Instruction printer}

\begin{itemize}
    \item That's it!
    \item Directives and labels handled for us
    \begin{itemize}
        \item Can emit target-specific syntax if we wish
    \end{itemize}
\end{itemize}

\examplebox[firstline=1,lastline=10]{ex1/ex1.s}

\end{frame}

%%%%%%%%%%%%%%%%%%%%%%%%%%%%%%%%%%%%%%%%%%%%%%%%%%%%%%%%%%%%%%%%%%%%%%%%%%%%%%%%

\talksection{Instruction encoding}

\begin{frame}{Instruction encoding}

\begin{itemize}
    \item A few new classes:
    \begin{itemize}
        \item An MCObjectStreamer (again, stock)
        \item LEGMCCodeEmitter
        \item LEGObjectWriter
        \item LEGAsmBackend
    \end{itemize}
    \item You will also need your LEGAsmPrinter
\end{itemize}

\end{frame}

%%%%%%%%%%%%%%%%%%%%%%%%%%%%%%%%%%%%%%%%%%%%%%%%%%%%%%%%%%%%%%%%%%%%%%%%%%%%%%%%

\begin{frame}[fragile]{Instruction encoding}

Example encoding:

\vspace{1ex}

\begin{tabular}{|c|c|c|c|c|c|c|c|}
\hline 
31...28 & 27...25 & 24...21 & 20 & 19...16 & 15...12 & 11...4 & 3...0\tabularnewline
\hline 
1110 & 000 & opcode & 0 & src1 & dst & 00000000 & src2\tabularnewline
\hline 
\end{tabular}

\vspace{2ex}

How can we achieve this?

\vspace{1ex}

\begin{codebox}[commandchars=\\\{\}]
def ADDrr : InstLEG<(outs GRRegs:$dst),
                    (ins GRRegs:$src1, GRRegs:$src2),
                    "add $dst, $src1, $src2",
                    [(set i32:$dst, (add i32:$src1, i32:$src2))]>;
\end{codebox}
\codecaption[-10.2ex]{LEGInstrInfo.td}

\end{frame}

%%%%%%%%%%%%%%%%%%%%%%%%%%%%%%%%%%%%%%%%%%%%%%%%%%%%%%%%%%%%%%%%%%%%%%%%%%%%%%%%

\begin{frame}[fragile]{Instruction encoding}

\begin{itemize}
    \item TableGen recognises the 'Inst' field:
\end{itemize}

\begin{codebox}
class InstLEG< ... > : Instruction {
  field bits<32> Inst;
  ...
}
\end{codebox}
\codecaption{LEGInstrFormats.td}

\begin{itemize}
    \item Used to define the binary encoding of each instruction in TableGen:
\end{itemize}

\begin{codebox}
def ADDrr : InstLEG< ... > {
  let Inst{31-25} = 0b110000;
  let Inst{24-21} = 0b1100;      // Opcode
  let Inst{20}    = 0b0;
  let Inst{11-4}  = 0b00000000;
}
\end{codebox}
\codecaption{LEGInstrInfo.td}

\end{frame}

%%%%%%%%%%%%%%%%%%%%%%%%%%%%%%%%%%%%%%%%%%%%%%%%%%%%%%%%%%%%%%%%%%%%%%%%%%%%%%%%

\begin{frame}[fragile]{Instruction encoding}

\begin{itemize}
    \item For operand-based encoding, need bit fields with the same names as the operands:
\end{itemize}

\begin{codebox}[commandchars=\\\[\]]
def ADDrr : InstLEG<(outs GRRegs:\codeemphc[$dst]),
                    (ins GRRegs:\codeempha[$src1], GRRegs:\codeemphb[$src2]) ... > {
  bits<4> \codeempha[src1]; bits<4> \codeemphb[src2]; bits<4> \codeemphc[dst];
  let Inst{31-25} = 0b110000;
  let Inst{24-21} = 0b1100;      // Opcode
  let Inst{20}    = 0b0;
  let Inst{19-16} = \codeempha[src1];        // Operand 1
  let Inst{15-12} = \codeemphc[dst];         // Destination
  let Inst{11-4}  = 0b00000000;
  let Inst{3-0}   = \codeemphb[src2];        // Operand 2
\end{codebox}
\codecaption{LEGInstrInfo.td}

\begin{itemize}
    \item \texttt{LEGMCCodeEmitter::getMachineOpValue()} will be called on each operand
\end{itemize}

\end{frame}

%%%%%%%%%%%%%%%%%%%%%%%%%%%%%%%%%%%%%%%%%%%%%%%%%%%%%%%%%%%%%%%%%%%%%%%%%%%%%%%%

\begin{frame}[fragile]{Instruction encoding}

\begin{itemize}
  \item \texttt{LEGMCCodeEmitter::getMachineOpValue()} returns the binary encoding of a given operand...
\end{itemize}

\begin{codebox}
unsigned LEGMCCodeEmitter::
getMachineOpValue(const MCInst &MI, const MCOperand MO,
                  SmallVectorImpl<MCFixup> &Fixups,
                  const MCSubtargetInfo &STI) const {
  if (MO.isReg()) {
    return
      CTX.getRegisterInfo()->getEncodingValue(MO.getReg());
  } if (MO.isImm()) {
    return static_cast<unsigned>(MO.getImm());
  }
  /* ... */
}
\end{codebox}
\codecaption{LEGMCCodeEmitter.cpp}

\begin{itemize}
  \item ...which is placed, masked, and shifted into position by TableGen-erated code
\end{itemize}

\end{frame}

%%%%%%%%%%%%%%%%%%%%%%%%%%%%%%%%%%%%%%%%%%%%%%%%%%%%%%%%%%%%%%%%%%%%%%%%%%%%%%%%

\begin{frame}[fragile]{Relocations and fixups}

\begin{itemize}
    \item For values that need fixing up, record the relocation and return zero
    \item LLVM will keep track of the relocation for us and help us fix it up later
\end{itemize}

\begin{codebox}
unsigned LEGMCCodeEmitter::
getMachineOpValue(const MCInst &MI, const MCOperand MO,
                  SmallVectorImpl<MCFixup> &Fixups,
                  const MCSubtargetInfo &STI) const {
  /* ... */

  assert(MO.isExpr()); // MO must be an expression
  
  const MCExpr *Expr = MO.getExpr();
  const MCExpr::ExprKind Kind = Expr->getFixupKind();

  Fixups.push_back(MCFixup::Create(0, Expr, Kind));
  return 0;
}
\end{codebox}
\codecaption{LEGMCCodeEmitter.cpp}

\end{frame}

%%%%%%%%%%%%%%%%%%%%%%%%%%%%%%%%%%%%%%%%%%%%%%%%%%%%%%%%%%%%%%%%%%%%%%%%%%%%%%%%

\begin{frame}[fragile]{Relocations and fixups}

\begin{itemize}
    \item Defining a target-specific fixup:
\end{itemize}

\begin{codebox}
enum Fixups {
  fixup_leg_movt_hi16_pcrel = FirstTargetFixupKind,
  fixup_leg_movw_lo16_pcrel,
  
  LastTargetFixupKind,
  NumTargetFixupKinds = LastTargetFixupKind - FirstTargetFixupKind
};
\end{codebox}
\codecaption{LEGFixups.h}

\begin{codebox}
const MCFixupKindInfo& getFixupKindInfo(MCFixupKind K) const {
  const static MCFixupKindInfo I[LEG::NumTargetFixupKinds] = {
    // Name          Offset Size Flags
    { "fixup_leg_movt_hi16_pcrel", 0,  32, MCFixupKindInfo::FKF_IsPCRel },
    { "fixup_leg_movw_lo16_pcrel", 0,  32, MCFixupKindInfo::FKF_IsPCRel },
  };
  /* ... */
  return I[K - FirstTargetFixupKind];
}
\end{codebox}
\codecaption{LEGAsmBackend.cpp}

\end{frame}

%%%%%%%%%%%%%%%%%%%%%%%%%%%%%%%%%%%%%%%%%%%%%%%%%%%%%%%%%%%%%%%%%%%%%%%%%%%%%%%%

\begin{frame}{Relocations and fixups}

\begin{itemize}
    \item We must then implement some hooks
    \item These are called at the end once the section layouts have been finalized
    \item \texttt{LEGAsmBackend::processFixupValue()}
    \begin{itemize}
        \item Adjusts the fixup value, e.g., splitting the value across non-contiguous fields
    \end{itemize}
    \item \texttt{LEGAsmBackend::applyFixup()}
    \begin{itemize}
        \item Patches the fixed-up value into the binary stream
    \end{itemize}
\end{itemize}

\end{frame}

%%%%%%%%%%%%%%%%%%%%%%%%%%%%%%%%%%%%%%%%%%%%%%%%%%%%%%%%%%%%%%%%%%%%%%%%%%%%%%%%

\talksection{Selection DAG manipulation}

\begin{frame}[fragile]{Custom SelectionDAG nodes}

\begin{itemize}
    \item To represent target-specific operations in the DAG
    \begin{itemize}
        \item Example: 32-bit immediate move
    \end{itemize}
    \item How?
    \begin{itemize}
        \item Add a value in the \texttt{LEGISD} enum
        \item Update \texttt{LEGTargetLowering::getTargetNodeName()}
        \item Add TableGen node definitions
        \begin{itemize}
            \item Type definition: number of inputs, outputs, constraints
            \item Node definition: tablegen name, opcode, type
        \end{itemize}
    \end{itemize}
    \item Custom nodes can be used in TableGen selection patterns
\end{itemize}

\begin{codebox}
def MoveImm32Ty : SDTypeProfile<1, 1, [
  SDTCisSameAs<0, 1>, SDTCisInt<0>
]>;

def movei32 : SDNode<"LEGISD::MOVi32", MoveImm32Ty>;
\end{codebox}
\codecaption{LEGOperators.td}

\end{frame}

%%%%%%%%%%%%%%%%%%%%%%%%%%%%%%%%%%%%%%%%%%%%%%%%%%%%%%%%%%%%%%%%%%%%%%%%%%%%%%%%

\begin{frame}[fragile]{Custom DAG lowering}

\begin{itemize}
    \item To handle DAG nodes in a special way
    \begin{itemize}
        \item Replaces an existing node with one or more other DAG nodes
        \item Matches nodes by opcode (e.g. \texttt{ISD::Constant})
        \item Matches nodes by type (e.g. \texttt{i32})
    \end{itemize}
    \item How?
    \begin{itemize}
        \item Call \texttt{setOperationAction(\emph{nodeOpcode}, type, Custom)}
        \item Create a function to handle it (e.g. \texttt{LowerOPCODE})
        \item Update \texttt{LowerOperation} to call \texttt{LowerOPCODE}
    \end{itemize}
    \item This all hapens in \texttt{LEGTargetLowering} (\texttt{LEGISelLowering.cpp})
\end{itemize}

\end{frame}

%%%%%%%%%%%%%%%%%%%%%%%%%%%%%%%%%%%%%%%%%%%%%%%%%%%%%%%%%%%%%%%%%%%%%%%%%%%%%%%%

\begin{frame}[fragile]{Custom DAG lowering: \texttt{LowerOPCODE}}

\begin{itemize}
    \item \texttt{LowerOPCODE} takes a DAG node (\texttt{Op}) and returns a DAG node
    \item You can:
    \begin{itemize}
        \item Return a different node
        \item Return \texttt{Op} → no change
        \item Return \texttt{SDValue()} → node not supported (LLVM will expand it)
    \end{itemize}
    \item All of the above can be done conditionally (e.g. depending on \texttt{VT})
\end{itemize}

\begin{codebox}
SDValue LEGTargetLowering::LowerConstant(SDValue Op,
                                         SelectionDAG &DAG) const {
  EVT VT = Op.getValueType();
  ConstantSDNode *Val = cast<ConstantSDNode>(Op.getNode());
  SDValue TargetVal = DAG.getTargetConstant(Val->getZExtVaue(),
                                            MVT::i32);
  return DAG.getNode(LEGISD::MOVi32, VT, TargetVal);
}
\end{codebox}
\codecaption{LEGISelLowering.cpp}

\end{frame}

%%%%%%%%%%%%%%%%%%%%%%%%%%%%%%%%%%%%%%%%%%%%%%%%%%%%%%%%%%%%%%%%%%%%%%%%%%%%%%%%

\begin{frame}[fragile]{Creating SelectionDAG nodes}

\begin{itemize}
    \item Simply call \texttt{DAG.getNode()} with these arguments:
    \begin{itemize}
        \item Node opcode (e.g. \texttt{LEGISD::MOVi32}), type, operand(s)
    \end{itemize}
    \item Nodes are target-independent (\texttt{ISD}) or not (\texttt{LEGISD})
    \item Use \texttt{DAG.getMachineNode()} in LEGISelDAGToDAG
\end{itemize}

\begin{codebox}
SDValue LEGTargetLowering::LowerConstant(SDValue Op,
                                         SelectionDAG &DAG) const {
  EVT VT = Op.getValueType();
  ConstantSDNode *Val = cast<ConstantSDNode>(Op.getNode());
  SDValue TargetVal = DAG.getTargetConstant(Val->getZExtVaue(),
                                            MVT::i32);
  return DAG.getNode(LEGISD::MOVi32, VT, TargetVal);
}
\end{codebox}
\codecaption{LEGISelLowering.cpp}

\end{frame}

%%%%%%%%%%%%%%%%%%%%%%%%%%%%%%%%%%%%%%%%%%%%%%%%%%%%%%%%%%%%%%%%%%%%%%%%%%%%%%%%

\begin{frame}{Lowering to multiple instructions}

\begin{minipage}[t]{0.58\linewidth}
    \begin{itemize}
        \item Example: 32-bit immediate load
        \begin{itemize}
            \item MOVW: loads 16-bit 'low' part
            \item MOVT: loads 16-bit 'high' part
        \end{itemize}
        %\item The final value is defined by the second instruction
        
        \item The two instructions must be ordered
        \begin{itemize}
            \item \texttt{MOVW} clears the 'high' part
            \item \texttt{MOVT, MOVW} gives the wrong result
            \item Make \texttt{MOVT} read the output of \texttt{MOVW}
        \end{itemize}
        \item Example: \texttt{0x00010002}
        
        %\item Different ways to implement this:
        %\begin{itemize}
        %    \item At the DAG level, in LEGISelDAGToDAG.cpp
        %    \item Using a pseudo-instruction (placeholder)
        %\end{itemize}
    \end{itemize}
\end{minipage}
\begin{minipage}[t]{0.41\linewidth}
    \begin{figure}
        \vspace{-3.5ex}
        \includegraphics[width = 1.00\textwidth]{examples/ex5/ex5-post-isel.pdf}
    \end{figure}
\end{minipage}

\end{frame}

%%%%%%%%%%%%%%%%%%%%%%%%%%%%%%%%%%%%%%%%%%%%%%%%%%%%%%%%%%%%%%%%%%%%%%%%%%%%%%%%

\begin{frame}[fragile]{Lowering to multiple instructions}

\begin{itemize}
    \item To define MOVT, we need an extra operand ('\texttt{fakesrc}')
    \begin{itemize}
        \item Source and destination registers must be the same
        \item Use '\texttt{Constraints}' in Tablegen
    \end{itemize}
\end{itemize}

\begin{codebox}[commandchars=\\¬|]
def MOVTi16 : InstLEG<(outs GRRegs:$dst),
                      (ins \codeempha¬GRRegs:$fakesrc|, i32imm:$src),
                      "movw $dst, $src",
                      [/* No pattern */]> {
  let Constraints = "$src1 = $dst";
}
\end{codebox}
\codecaption{LEGInstrInfo.td}

\begin{itemize}
    \item Different ways to emit multiple instructions from one DAG node
    \begin{itemize}
        \item Using custom C++ instruction selection code
            \begin{itemize}
                \item Not covered here
            \end{itemize}
        \item Using a pseudo-instruction as a placeholder
    \end{itemize}
\end{itemize}

\end{frame}

%%%%%%%%%%%%%%%%%%%%%%%%%%%%%%%%%%%%%%%%%%%%%%%%%%%%%%%%%%%%%%%%%%%%%%%%%%%%%%%%

\begin{frame}[fragile]{Lowering to multiple instructions}

\begin{itemize}
    \item Using a pseudo instruction
    \begin{itemize}
        \item Behaves like a placeholder for 'real' machine instruction(s)
        \item Lowered by a target hook into these instruction instructions
        \item Can be selected from the custom DAG node we previously defined
    \end{itemize}
\end{itemize}

\begin{codebox}
def MOVi32 : InstLEG<(outs GRRegs:$dst), (ins i32imm:$src), "",
                     [(set i32:$dst, (movei32 i32imm:$src))]> {
  let isPseudo = 1;
}
\end{codebox}
\codecaption{LEGInstrInfo.td}

\end{frame}

%%%%%%%%%%%%%%%%%%%%%%%%%%%%%%%%%%%%%%%%%%%%%%%%%%%%%%%%%%%%%%%%%%%%%%%%%%%%%%%%

\begin{frame}[fragile]{Lowering to multiple instructions}

\begin{itemize}
    \item The pseudo is lowered by a target hook
\end{itemize}


\begin{codebox}
bool LEGInstrInfo::expandPostRAPseudo(MachineInstr *MI) {
   if (MI->getOpcode() != LEG::MOVi32) return false;
   BuildMI(..., get(LEG::MOVWi16), Dst).addImm(Lo);
   BuildMI(..., get(LEG::MOVTi16), Dst).addReg(Dst).addImm(Hi);
   MBB.erase(MI);
   return true;
}
\end{codebox}
\codecaption{LEGInstrInfo.cpp}

\end{frame}

%%%%%%%%%%%%%%%%%%%%%%%%%%%%%%%%%%%%%%%%%%%%%%%%%%%%%%%%%%%%%%%%%%%%%%%%%%%%%%%%

\begin{frame}{Lowering to multiple instructions}

Example IR:
\examplebox[firstline=6,lastline=9]{ex5/ex5.ll}

Resulting assembly:
\examplebox[firstline=7,lastline=11,gobble=1]{ex5/ex5.s}

\end{frame}
